\subsection{TF-IDF logistička regresija kao podsustav klasifikacije}

\paragraph{Motivacija.}
Ovaj podsustav promatra komentar kao tekstualni niz i traži lokalne leksičke i znakovne obrasce. Za razliku od semantičkog SVM-a, ovdje se ne koristi gusto semantičko ugrađivanje rečenice; za razliku od stablastog modela, ne koriste se ručno zadane strukturne značajke. Model izravno uči iz riječi, parova riječi i kratkih nizova znakova. Takav pristup je osobito koristan jer AI-generirani komentari često ponavljaju karakteristične fraze, prijelaze, ritam izražavanja, interpunkcijske obrasce i dijelove formulacija koji se mogu uhvatiti bez dubljeg razumijevanja značenja.

Za komentar $t$ ovaj podsustav daje vjerojatnost

\[
p_{\mathrm{tf}}(t) \approx P(y=1\mid t),
\]

gdje $y=1$ označava razred \texttt{AI}. Ta se vrijednost dalje koristi kao treći bazni ulaz završnog stacking modela.

\paragraph{TF-IDF prikaz teksta.}
Tekst se najprije pretvara u rijedak vektor značajki pomoću dvaju TF-IDF prikaza. Prvi prikaz koristi riječi i parove uzastopnih riječi, odnosno unigram i bigram značajke. Drugi prikaz koristi znakovne $n$-grame duljina od $3$ do $5$. U provedbi su to dvije komponente \texttt{FeatureUnion}-a: \texttt{word\_tfidf} i \texttt{char\_tfidf}. U oba se prikaza odbacuju značajke koje se pojavljuju u manje od $2$ komentara ili u više od $95\%$ komentara.

Za značajku $g$ i komentar $t$ neka je $c_{t,g}$ broj pojavljivanja značajke $g$ u komentaru. Budući da je u modelu uključeno sublinearno skaliranje frekvencije, lokalni dio težine glasi

\[
\operatorname{tf}(t,g)=
\begin{cases}
1+\log c_{t,g}, & c_{t,g}>0,\\
0, & c_{t,g}=0.
\end{cases}
\]

Ako je $N$ broj komentara u skupu za učenje, a $\operatorname{df}(g)$ broj komentara u kojima se značajka $g$ pojavljuje, inverzna dokumentna frekvencija u standardnoj zaglađenoj inačici može se zapisati kao

\[
\operatorname{idf}(g)=\log\frac{1+N}{1+\operatorname{df}(g)}+1.
\]

Neobrađena TF-IDF težina zatim je

\[
\tilde x_g(t)=\operatorname{tf}(t,g)\operatorname{idf}(g),
\]

a dobiveni vektor se L2-normalizira. Time vrlo česte riječi dobivaju manju relativnu važnost, dok izrazi koji su dovoljno česti da budu pouzdani, ali nisu svugdje prisutni, dobivaju jači signal.

\paragraph{Riječni i znakovni $n$-grami.}
Riječni dio modela koristi $n$-grame duljina $1$ i $2$. On hvata izraze poput pojedinačnih ključnih riječi, čestih fraza i lokalnih kombinacija riječi. Znakovni dio koristi $n$-grame duljina od $3$ do $5$, pa može uhvatiti i obrasce unutar riječi, dijelove sufiksa, interpunkcijske prijelaze, kontrakcije i tipične fragmente koji ne moraju biti cijele riječi.

Konačni vektor dobiva se konkatenacijom ta dva prikaza:

\[
x(t)=
\begin{bmatrix}
x_{\mathrm{word}}(t)\\
x_{\mathrm{char}}(t)
\end{bmatrix}
\in\mathbb{R}^{d}.
\]

U finalnom modelu spremljenom za stacking rječnik riječnih TF-IDF značajki ima $793698$ elemenata, a rječnik znakovnih TF-IDF značajki $763983$ elemenata. Ukupno je dakle

\[
d=1557681
\]

rijetkih značajki. Iako je dimenzija vrlo velika, svaki pojedini komentar aktivira samo mali dio tih značajki, pa je računanje učinkovito.

\paragraph{Logistička regresija nad rijetkim tekstualnim vektorom.}
Nakon vektorizacije trenira se logistička regresija s L2-regularizacijom i uravnoteženim težinama razreda. U provedbi se koristi \texttt{solver=liblinear} i \texttt{max\_iter=2000}. Za razred $c$ model računa linearni rezultat

\[
s_c(x)=w_c^\top x+b_c.
\]

Vjerojatnost razreda dobiva se normalizacijom tih rezultata. Za AI razred to možemo zapisati kao

\[
p_{\mathrm{tf}}(t)
=P(y=\texttt{AI}\mid x(t))
=\frac{\exp(s_{\texttt{AI}}(x(t)))}{\exp(s_{\texttt{AI}}(x(t)))+\exp(s_{\texttt{HUMAN}}(x(t)))}.
\]

U binarnom slučaju taj je zapis ekvivalentan sigmoidnoj funkciji nad jednom linearnom kombinacijom značajki. Važno je da model ostaje linearan u TF-IDF prostoru: pojedini koeficijent govori koliko prisutnost odgovarajućeg $n$-grama pomiče odluku prema jednom ili drugom razredu.

Optimizacija se može opisati minimizacijom regulariziranog logističkog gubitka

\[
\min_w\;
\sum_{i=1}^{n}\alpha_{y_i}
\left[-y_i\log p_i-(1-y_i)\log(1-p_i)\right]
+\lambda\|w\|_2^2,
\]

gdje su $\alpha_{y_i}$ težine razreda. U provedbi se koristi \texttt{class\_weight=balanced}, pa rjeđi razred dobiva veću težinu u gubitku. Time se smanjuje opasnost da model jednostavno favorizira brojniji razred.

\paragraph{Zašto je ovakav model prikladan.}
TF-IDF logistička regresija vrlo je jednostavna, ali je posebno jaka za zadatke u kojima razredi imaju prepoznatljive tekstualne obrasce. Velika prednost je što model ne mora učiti duboki jezični prikaz od početka: on samo traži koji su $n$-grami statistički povezani s razredom. Zbog toga je brz, jeftin za treniranje i vrlo dobar kao bazni model.

Druga važna prednost je komplementarnost. Semantički SVM može uhvatiti značenje komentara, a stablasti model makro-strukturu, ali TF-IDF logistička regresija hvata sitne lokalne tragove koje ta dva modela mogu propustiti. Upravo zato je korisna kao treći ulaz završnog meta-modela.

\paragraph{Izlaz podsustava.}
Spremljeni \texttt{sklearn} pipeline sadrži i TF-IDF vektorizatore i logistički klasifikator. Zato se novi komentar može izravno predati modelu, a funkcija \texttt{predict\_proba} vraća vjerojatnosti razreda. U finalnom spremljenom modelu poredak klasa je

\[
\texttt{AI},\quad \texttt{HUMAN},
\]

pa se za izlaz podsustava uzima stupac koji pripada razredu \texttt{AI}. Kad bi se model koristio samostalno, binarna odluka mogla bi se zadati pravilom

\[
\hat y = \mathbf{1}\{p_{\mathrm{tf}}(t)\ge 0.5\}.
\]

U završnom sustavu ponovno je važnija sama vjerojatnost $p_{\mathrm{tf}}(t)$, jer je ona ulaz u stacking meta-model.

\paragraph{Rezultati na razvojnom skupu.}
Na razvojnom skupu finalni TF-IDF logistički model postiže točnost $0.98085$, preciznost $0.98627$, odziv $0.97447$, mjeru $F_1=0.98033$ i pokazatelj $\mathrm{ROC\ AUC}=0.99780$. Budući da model vraća vjerojatnosti, korisne su i vjerojatnosne metrike: logaritamski gubitak iznosi $0.06739$, a Brierova mjera $0.01626$.

Konfuzijska matrica na razvojnom skupu iznosi

\[
\begin{array}{c|cc}
 & \hat y=0 & \hat y=1 \\
\hline
y=0 & 21290 & 281 \\
y=1 & 529 & 20189
\end{array}
\]

Ovaj podsustav je najjači pojedinačni bazni model u trenutnoj izvedbi. Ipak, završni meta-model dodatno poboljšava rezultat jer iskorištava slučajeve u kojima se TF-IDF signal može potvrditi ili ispraviti semantičkim i strukturnim signalom.

\paragraph{Mogućnost jednostavnog poboljšanja.}
Najjednostavnije poboljšanje ovoga podsustava bilo bi sustavno podešavanje hiperparametara TF-IDF prikaza i logističke regresije: primjerice vrijednosti \texttt{min\_df}, \texttt{max\_df}, raspona riječnih i znakovnih $n$-grama te regularizacijskog parametra $C$. Drugi prirodan korak je dodatna kalibracija vjerojatnosti, jer je završnom stacking modelu važna kvaliteta samih vjerojatnosti, a ne samo točnost binarne odluke. Takve izmjene ne bi promijenile osnovnu ideju modela, nego bi samo preciznije prilagodile istu linearnu metodu podacima.